\section{R ile korelasyon uygulaması}
\label{sec:r-correlation}

\textbf{Soru.} Bölümdeki 20 öğrencinin uyku süresi ve kaygı puanları için
saçılım grafiğini, Pearson ve Spearman korelasyonlarını ve Pearson güven
aralığını üretin.

\begin{lstlisting}[style=prismR]
uyku <- c(4.5,5.0,5.2,5.5,5.8,6.0,6.1,6.3,6.5,6.7,
          6.8,7.0,7.1,7.3,7.5,7.6,7.8,8.0,8.2,8.5)
kaygi <- c(68,52,75,60,66,48,70,55,64,58,
           45,62,53,46,60,43,55,49,57,40)

plot(uyku, kaygi, pch = 19,
     xlab = "Uyku suresi", ylab = "Kaygi puani")
abline(lm(kaygi ~ uyku), col = "blue", lwd = 2)

pearson <- cor.test(uyku, kaygi, method = "pearson")
spearman <- cor.test(uyku, kaygi, method = "spearman",
                     exact = FALSE)
pearson
spearman

# Tek gozlemin duyarlilik etkisi
cor(uyku[-3], kaygi[-3], method = "pearson")
\end{lstlisting}

\begin{outputguidebox}
\textbf{Çözüm ve kontrol değerleri.} Pearson korelasyonu
$r=-0{,}5849$, $t(18)=-3{,}061$, $p=0{,}00675$'tir. Spearman katsayısı
yaklaşık $\rho_s=-0{,}5493$, $p=0{,}0121$ olur. Kesin güven aralığı
\texttt{cor.test} çıktısından okunur. Katsayıların yönü aynıdır; sayısal
büyüklüklerinin farklı olması bir yazılım hatası değildir.
\end{outputguidebox}

\textbf{Yorum.} Pearson doğrusal birlikte değişimi, Spearman sıraların tekdüze
ilişkisini özetler. İki katsayıyı sonuçtan sonra daha küçük $p$-değerini seçmek
için değil, veri yapısı ve araştırma hedefi doğrultusunda kullanın. Gözlemsel
kesitsel veri nedensel uyku etkisini göstermez.

\textbf{Pekiştirme ve yanıt.} Üçüncü gözlemi çıkarmak yalnız duyarlılık
karşılaştırmasıdır. Kayıt hatası veya önceden belirlenmiş bir dışlama gerekçesi
yoksa gözlem sırf katsayı değiştiği için silinmemelidir.